\hwhat{The question is whether the kickoff text is the target into which a goal change quenches the swarm. Each day-1 content centroid is scored against its own kickoff and 32 others (genericness-corrected swap null), in two embedding models. The centroid picks out its own kickoff (median percentile 1.0; top-1 18/33 with bge, 20/33 with gte). It fails only where the kickoff names no shared target. In \#51 each agent lands on its own private goal (0.95 / 0.96), and a reassignment moves one agent within a day (+0.61 / +0.56). Round 2 asks how text acts after day 1. A human message pulls each reader at its first message after the read-out call ($\Delta\approx0.09$). Reading adds $+0.11$ over unread in-flight messages at matched lag. Agent-written plans do not pull readers, and two offered targets do not split the swarm. The reserved confirmation data are not used yet.}

\hmodel{Vector spins with a quench (model 11). Each agent's whitened statement vector $\mathbf s_i$ relaxes toward a target $\hat{\mathbf t}$ with depth $\lambda$. It carries an offset $\mathbf a_i$ (style, family, old project), a herding pull $J\mathbf m$ and noise. H54 says the kickoff text gives the target: $\hat{\mathbf t}\approx\hat{\mathbf k}$, the kickoff embedding. In the Potts version (model 10), the kickoff is a field on the projects it names, which freeze at once. Rivals: no target, inertia, genre, agent plans, family fields.}

\hmath{\vspace{-8pt}
\begin{gather*}
\mathbf s_i(\text{day }1)\propto\lambda\,\hat{\mathbf t}+\mathbf a_i+J\,\mathbf m+\boldsymbol\eta_i,\qquad \hat{\mathbf t}\approx\hat{\mathbf k}_p\\
\hat S_{pq}=\cos(\hat{\mathbf m}_p,\hat{\mathbf k}_q)-\langle\cdot\rangle_{p'\ne q},\qquad \pi_p=\Pr_{q\ne p}\big[\hat S_{pq}<\hat S_{pp}\big]
\end{gather*}\vspace{-8pt}}

\hobs{../figures/summary_obs_page.pdf}{(a) Own-kickoff percentile of each day-1 centroid among 33 kickoffs. Squares mark an own kickoff ranked first. Grey: the 90\% band of the no-target null. Open circles: kickoffs that name no shared target (free weeks, \#44, \#51); all four failures are among them. (b) Kickoff-frozen H31 projects against all others: the share whose name matches the goal text, and the share that pre-existed (the carry-over rival). Frozen projects carry the goal's name and are new.}

\htelos{Write the target into the kickoff text. Day-1 talk lands on it in both embedding models, and the frozen project carries the goal's words (9/13). A private goal text pins each agent to its own target for weeks. A human message moves the agents who read it within their next message, but the pull is gone within about 10 minutes, and half of it shows up in unread messages too (convergence). Do not expect an agent's plan to steer the others: readers move away from it. Offering two tasks in one kickoff gives one mixed swarm, not two teams; assign people explicitly if you want a split. Text specificity is not a validated dial.}

\hbottom{The kickoff or private goal text is the day-1 quench target in both embedding models; later human messages re-quench readers briefly through reading, agent plans do not, and specificity does not set the spread.}

\hresults{{\footnotesize\noindent\begin{tabular}{@{}p{0.34\columnwidth}p{0.45\columnwidth}p{0.15\columnwidth}@{}}
\textbf{Prediction} & \textbf{Observed (bge / gte)} & \textbf{Verdict}\\
P1 own kickoff identified & median $\pi$ 1.0; top-1 18/33 / 20/33; $\rho_{\rm models}=0.75$ & supported\\
P1c move points at kickoff & 0.91 / 0.97 & supported\\
P2 frozen = named & 9/13; goal-text enrichment 3.2 & mixed\\
P3 specificity sets spread & $\rho=+0.17$ / $+0.19$ & failed\\
\#51 own goal; NE38 & 0.95 / 0.96; DiD $+0.61$ / $+0.56$ & supported\\
R3 read-out re-quench & $\Delta$ 0.090 / 0.088, 181 msgs & supported\\
R3 read $-$ in flight (matched) & $+0.11\,[0.05,0.17]$ / $+0.12$ & supported\\
R2 first-plan readers & $-0.10$ / $-0.10$; read$-$in flight unpowered & failed\\
R1 \#12 motion; team domains & percentile 1.0; $Q=+0.002$ ($p$ 0.40) & mixed\\
R1 \#19/\#21 spontaneous domains & none (pct $\le0.87$) & as predicted\\
HH184 family susceptibility & lab $p=0.66$ / $0.67$ & failed\\
\end{tabular}}}

\hobsb{../figures/round2.pdf}{(a) Round 2, R3: pull of a human message on each reader, against the lag of the reader's first message after it (181 messages). Lines: messages produced after the read-out call. Orange: unread messages produced in flight (median lag 8 s). At matched lag, reading roughly doubles the pull. (b) Own-kickoff percentile per period in the two embedding models; labels mark periods below 0.75 in either.}

\hcaveats{The target is a text-embedding proxy, and quoting is part of the effect (removing near-echoes lowers top-1 to 0.45 or 0.30). P1, \#51 and NE38 hold in both models; round-1's room-centroid re-quench, \#38's room swap and agent susceptibility do not agree across models. Only 93 in-flight pairs exist (83 at lag $<30$ s), so the read-out contrast rests on fast replies; a paired within-agent contrast is null because the read message comes later and the pull decays (post hoc). The plan test has no power for its causal contrast. ``Named'' is name-token matching. 33 kickoffs, 22 in regime I. \#23 is excluded to keep H10's confirmation blind.}

\hnext{Run the confirmatory test on reserved data (14 kickoffs, the \#51 tail and reserved human messages), with the read-out contrast added. Fit the read-out pulse decay on the per-call clock. Test LLM-rated specificity (R4) and the remanence mechanism (R6).}
